\section*{Chapitre \ref{ch:b3:atmospheric-dynamics} --- Dynamique atmosphérique}

\begin{solution}{exo:b3:atmospheric-dynamics:1}
(a)~$\zeta = 2\omega = 2\times2\pi/86\,400 = \qty{1.45e-4}{s^{-1}}$.
(b)~$f = \num{7.29e-5}$ et \qty{1.26e-4}{s^{-1}}~: le \omterm{def:b3:atmospheric-dynamics:eddy}{tourbillon} a environ
deux fois la \omterm{def:b3:atmospheric-dynamics:vorticity-kinds}{vorticité planétaire} à \ang{30}~N et un peu plus que celle à
\ang{60}~N.
\end{solution}

\begin{solution}{exo:b3:atmospheric-dynamics:2}
$\zeta = f_{30} - f_{40} = \num{7.29e-5} - \num{9.37e-5} =
\qty{-2.08e-5}{s^{-1}}$~: anticyclonique.
\end{solution}

\begin{solution}{exo:b3:atmospheric-dynamics:3}
$k = 2\pi/(5\times10^6) = \qty{1.257e-6}{m^{-1}}$; $c = -\beta/k^2 =
-\num{1.62e-11}/\num{1.58e-12} = \qty{-10.3}{m/s}$~: vers l'ouest.
\end{solution}

\begin{solution}{exo:b3:atmospheric-dynamics:4}
$\beta = 2\Omega\cos\ang{60}/a = \qty{1.14e-11}{m^{-1}.s^{-1}}$;
$\lambda_s = 2\pi\sqrt{10/\num{1.14e-11}} = \qty{5870}{km}$~; le cercle,
\qty{20015}{km}, contient $n_s = 3.4$ ondes.
\end{solution}

\begin{solution}{exo:b3:atmospheric-dynamics:5}
$q = f/h$ se conserve~: sur la chaîne, $\zeta = f(8/10) - f =
\qty{-2.06e-5}{s^{-1}}$, anticyclonique~; sur le bassin, $f(12/10) - f =
\qty{+2.06e-5}{s^{-1}}$, cyclonique. Là où les colonnes s'étirent, à l'aval
de la chaîne et sur le bassin, une vorticité cyclonique se crée~: c'est là que
les dépressions ont tendance à se former.
\end{solution}

\begin{solution}{exo:b3:atmospheric-dynamics:6}
$k = 5/(a\cos\ang{45}) = \qty{1.11e-6}{m^{-1}}$, $\beta/k^2 =
\qty{13.1}{m/s}$~: $c = \qty{1.9}{m/s}$, $c_g = \qty{28.1}{m/s}$. L'énergie
parcourt \qty{5000}{km} en \num{2.1}~jours, la figure d'onde en \num{31}~jours.
\end{solution}

\begin{solution}{exo:b3:atmospheric-dynamics:7}
(a)~$k^2 = \beta/U - l^2 = \num{1.079e-12} - \num{2.74e-13} =
\qty{8.05e-13}{m^{-2}}$, $k = \qty{8.97e-7}{m^{-1}}$~: $\lambda =
\qty{7000}{km}$ et $n = 4.0$, moins d'ondes que pour $l = 0$. (b)~$l^2 =
\num{1.097e-12} > \beta/U = \num{1.079e-12}$~: $k^2$ serait négatif.
\end{solution}

\begin{solution}{exo:b3:atmospheric-dynamics:8}
$L_d = 0.012\times10^4/\num{1.031e-4} = \qty{1164}{km}$; $\sigma =
0.31\times\num{1.031e-4}\times\num{2e-3}/0.012 = \qty{5.3e-6}{s^{-1}}$, soit
un temps de repliement en $e$ de \num{2.2}~jours~; longueur d'onde $2\pi\times1164/1.61 =
\qty{4540}{km}$.
\end{solution}

\begin{solution}{exo:b3:atmospheric-dynamics:9}
$[v^*T^*] = (VT_0/2)\cos\delta$~: \qty{25}{K.m/s}, \qty{12.5}{K.m/s} et 0~;
flux de chaleur \num{25100}, \num{12600} et \qty{0}{W/m^2}. Quand la
température est décalée d'un quart de longueur d'onde par rapport au vent, autant
d'air chaud part vers le sud que vers le nord~: pas de transport.
\end{solution}

\begin{solution}{exo:b3:atmospheric-dynamics:10}
En linéarisant autour de $U$ avec $\zeta' = \nabla^2\psi'$ et $v' =
\partial_x\psi'$~: $(\partial_t + U\partial_x)\nabla^2\psi' + \beta\partial_x
\psi' = 0$. Avec $\psi' \propto \mathrm e^{\mathrm i(kx + ly - \omega t)}$,
$(\omega - Uk)K^2 = -\beta k$, donc $\omega = Uk - \beta k/K^2$. Par rapport à
l'air, la vitesse de phase vaut $c - U = -\beta/K^2 < 0$ pour toute onde.
\end{solution}

\begin{solution}{exo:b3:atmospheric-dynamics:11}
$\varphi_H = [5\times9.81\times12\,000\times0.25/(3\times(\num{7.292e-5})^2
\times(\num{6.371e6})^2)]^{1/2} = \qty{0.477}{rad}$, \ang{27.3}.
$u_M = \Omega a\varphi_H^2 = \qty{106}{m/s}$~; exactement, $\Omega a\sin^2
\varphi_H/\cos\varphi_H = \qty{110}{m/s}$~: deux fois et demie le
courant-jet observé.
\end{solution}

\begin{solution}{exo:b3:atmospheric-dynamics:12}
(a)~À travers la bande, l'interface est déplacée de $\eta = sy$, $-L/2 \le y
\le L/2$, par rapport à sa position horizontale. Une couche à contraste de densité réduit
déplacée de $\eta$ emmagasine $\tfrac12\rho g'\eta^2$ par unité de surface~; moyennée
sur la bande, $\langle\eta^2\rangle = s^2L^2/12$. (b)~$g' = 0.327$~:
\num{1.5e5} et \qty{6.1e5}{J/m^2}, contre $\tfrac12\times0.7\times15^2
\times10^4 = \qty{7.9e5}{J/m^2}$ d'énergie cinétique~: comparables.
\end{solution}

\begin{solution}{pb:b3:atmospheric-dynamics:1}
\textbf{1.} $f_{45} = \qty{1.031e-4}{s^{-1}}$, $f_{50} =
\qty{1.117e-4}{s^{-1}}$; $\zeta = f_{45} - f_{50} = \qty{-8.6e-6}{s^{-1}}$.

\textbf{2.} $R = V/|\zeta| = 15/\num{8.6e-6} = \qty{1750}{km}$~; sens horaire,
en revenant vers le sud.

\textbf{3.} $\zeta = f_{45}(9.5/10) - f_{50} = \qty{-1.38e-5}{s^{-1}}$, plus
anticyclonique.

\textbf{4.} $\zeta = f_{45} - f_{40} = \qty{+9.4e-6}{s^{-1}}$, cyclonique~: elle
revient vers le nord. Écartée dans un sens ou dans l'autre, la particule acquiert la vorticité
qui la ramène~: un effet de rappel, et une oscillation.

\textbf{5.} $\beta = \qty{1.62e-11}{m^{-1}.s^{-1}}$~; le cercle a une longueur de
\qty{28306}{km} long.

\textbf{6.} $n = 3$ à 8~: longueurs d'onde 9435, 7076, 5661, 4718, 4044 et
\qty{3538}{km}~; $c = -21.5$, $-5.5$, 1.9, 5.9, 8.3 et \qty{9.9}{m/s}~; $c_g
= 51.5$, 35.5, 28.1, 24.1, 21.7 et \qty{20.1}{m/s}.

\textbf{7.} $n = 3$ et 4 se déplacent vers l'ouest, $n \ge 5$ vers l'est. L'onde $n = 5$ met
$10^6/1.86 = \qty{5.4e5}{s}$, soit \num{6.2}~jours, à parcourir \qty{1000}{km}.

\textbf{8.} À l'aval, vers l'est~: à $c_g = \qty{24.1}{m/s}$, l'énergie parcourt
\qty{5000}{km} en \num{2.4}~jours.

\textbf{9.} $\lambda_s = 2\pi\sqrt{15/\num{1.62e-11}} = \qty{6050}{km}$,
$n_s = 4.7$.

\textbf{10.} $n_s = 5.7$ pour \qty{10}{m/s} et 4.1 pour \qty{20}{m/s}~: un
vent d'ouest plus fort favorise des ondes stationnaires plus longues, donc
moins nombreuses autour de l'hémisphère.

\textbf{11.} $n_s = 5.8$ à \ang{35}~N et 3.4 à \ang{55}~N.

\textbf{12.} $\zeta = f(8/10) - f = \qty{-2.06e-5}{s^{-1}}$~: la compression
au-dessus des montagnes rend la colonne anticyclonique, tandis que
l'étirement sous le vent la rend cyclonique. Les montagnes imposent au
courant une configuration fixe de vorticité, un forçage qui ne se déplace pas.

\textbf{13.} Un forçage fixe produit une réponse stationnaire~; la réponse est
la plus forte lorsque la configuration forcée est elle-même une onde libre
qui resterait immobile, c'est-à-dire la résonance à $\lambda_s$.

\textbf{14.} $\Lambda = 25/10^4 = \qty{2.5e-3}{s^{-1}}$~;
$|\partial T/\partial y| = \Lambda fT/g = \num{2.5e-3}\times\num{1.031e-4}
\times250/9.81 = \qty{6.6}{K}$ par \qty{1000}{km}.

\textbf{15.} $\mathrm{Ri} = 0.011^2/(\num{2.5e-3})^2 = 19$~: bien au-dessus de 0.25,
pas de turbulence de cisaillement.

\textbf{16.} $L_d = 0.011\times10^4/\num{1.031e-4} = \qty{1067}{km}$.

\textbf{17.} $2\pi\times1067/1.61 = \qty{4160}{km}$; $28306/4160 = 6.8$
ondes.

\textbf{18.} $2\pi\times1067/2.40 = \qty{2790}{km}$.

\textbf{19.} $\sigma = 0.31\times\num{1.031e-4}\times\num{2.5e-3}/0.011 =
\qty{7.3e-6}{s^{-1}}$, \qty{0.63}{jour^{-1}}.

\textbf{20.} $\exp(3\times0.627) = 6.6$.

\textbf{21.} La tempête, de \qty{4160}{km}, est plus courte que la longueur
d'onde stationnaire, \qty{6050}{km}~: en tant qu'onde de Rossby, elle est
entraînée vers l'est par le vent d'ouest, et les tempêtes se déplacent vers
l'est à travers les ondes longues plus lentes ou immobiles.

\textbf{22.} Vers l'ouest avec l'altitude~: ce n'est qu'à cette condition que
l'air qui se dirige vers le pôle est plus chaud que celui qui se dirige vers l'équateur, de
sorte que la chaleur est transportée vers le pôle et que l'\omterm{def:b3:atmospheric-dynamics:ape}{énergie potentielle disponible} est libérée.

\textbf{23.} La perturbation la plus rapide croît d'un facteur $e$ tous les
\textbf{\qty[text-series-to-math=true]{1.6}{jours}} (\num{1.59}~jours selon
le calcul).
\end{solution}
